\chapter*{CAPÍTULO 5. Validación del modelo SocialScrum mediante la técnica de juicio de expertos}
\addcontentsline{toc}{chapter}{CAPÍTULO 5. Validación del modelo SocialScrum mediante la técnica de juicio de expertos}
\setcounter{chapter}{5}
\setcounter{section}{0}
\setcounter{subsection}{0}
\setcounter{subsubsection}{0}
\renewcommand{\thesection}{\thechapter.\arabic{section}}
\renewcommand{\thesubsection}{\thesection.\arabic{subsection}}
\renewcommand{\thesubsubsection}{\thesubsection.\arabic{subsubsection}}

\setlength{\parindent}{0pt}

Este capítulo presenta la validación del modelo SocialScrum, con el fin de analizar su coherencia y alineación con los principios de la agilidad. La validación se aborda desde una perspectiva conceptual, considerando tanto la evaluación de la agilidad del modelo como la valoración experta de su diseño, previo a una eventual implementación empírica en contextos organizacionales reales.

Es importante señalar que SocialScrum prioriza deliberadamente los factores sociales y organizacionales del trabajo en equipo por encima de prácticas técnicas específicas de ingeniería de software. En consecuencia, los resultados de esta validación deben interpretarse como una validación conceptual del modelo desde la perspectiva ágil, quedando la evaluación de su desempeño técnico y operativo como una línea de trabajo futuro.
La estructura de este capítulo se organiza de la siguiente manera y se puede visualizar en la figura \ref{fig:estructura-capitulo5}:

\begin{enumerate}
    \item \textbf{Protocolo de validación mediante juicio de expertos}
    \item \textbf{Perfil y selección de expertos participantes}
    \item \textbf{Diseño del instrumento de validación}
    \item \textbf{Procedimiento de aplicación}
    \item \textbf{Análisis de resultados y retroalimentación}
\end{enumerate}

\section{Estrategia de validación del modelo SocialScrum}
La validación del modelo SocialScrum se fundamenta en un enfoque metodológico mixto que busca contrastar la viabilidad conceptual, la pertinencia operativa y el cumplimiento de los principios ágiles. Dado que el modelo propone una extensión socio-organizacional de Scrum, la estrategia de validación se divide en dos dimensiones complementarias: (i) una evaluación técnica de contenido mediante el juicio de expertos, y (ii) un análisis de alineación con el Manifiesto Ágil utilizando el framework AgilityPRO \cite{ortega2019agilityevaluation}.

Para garantizar que SocialScrum no solo sea una propuesta teórica, sino una herramienta de ingeniería ejecutable, el modelo se somete a la evaluación de cinco criterios de calidad fundamentales detallados en la Sección 5.3. Estos criterios permiten medir aspectos clave como la claridad, completitud, idoneidad, facilidad de uso y agilidad del modelo. La evaluación se realiza a través de un instrumento estructurado que combina preguntas cerradas con escala Likert y preguntas abiertas para capturar tanto valoraciones cuantitativas como cualitativas.

\subsection{Aplicación de AgilityPRO en la dimensión de Agilidad}

El uso de AgilityPRO (Agility in Processes) \cite{ortega2019agilityevaluation}, es crucial para determinar si SocialScrum mantiene la esencia de los valores y principios del Manifiesto Ágil a pesar de su enfoque en la deuda social. Este framework proporciona un soporte formal para evaluar la agilidad de los procesos mediante el análisis del cumplimiento de los principios ágiles, independientemente de la metodología base utilizada.
En esta validación, se utiliza el modelo de referencia AgilityRef para mapear cómo los elementos de SocialScrum (como la Social Sprint Retrospective o el Social Product Backlog) actúan como "aspectos de agilidad" que evidencian la implementación de uno o más principios ágiles. De este modo, la validación no depende de una percepción subjetiva, sino de un marco de referencia internacional que garantiza que el modelo apoya el desarrollo iterativo, la inspección frecuente y la centralidad en las personas.

\section{Criterios de evaluación considerados}
La validación del modelo SocialScrum se llevó a cabo mediante un cuestionario estructurado que incluyó tanto preguntas cuantitativas como cualitativas. El cuestionario se diseñó para evaluar diferentes aspectos del modelo, tales como su coherencia con los principios ágiles, su aplicabilidad práctica y su capacidad para mejorar la colaboración en equipos distribuidos.
\subsection{Criterios de evaluación}
Se definieron cinco criterios de evaluación clave para guiar la valoración del modelo SocialScrum por parte de los expertos:

\begingroup
\scriptsize
\renewcommand{\arraystretch}{1.4}
\setlength{\tabcolsep}{3pt}

\begin{longtable}{
    >{\RaggedRight\hspace{0pt}}p{3.0cm}   % Criterio
    >{\RaggedRight\hspace{0pt}}p{10.4cm}  % Descripción
    }

    % --- ENCABEZADO PRIMERA PÁGINA ---
    \caption{Criterios de evaluación y número de preguntas en el instrumento de validación}
    \label{tab:criterios-evaluacion}                                                                                                                                                                              \\
    \toprule
    \textbf{Criterio} & \textbf{Descripción}                                                                                                                                                                      \\
    \midrule
    \endfirsthead

    % --- ENCABEZADO PÁGINAS SIGUIENTES ---
    \multicolumn{2}{c}{{\bfseries \tablename\ \thetable{} -- Continuación}}                                                                                                                                       \\
    \toprule
    \textbf{Criterio} & \textbf{Descripción}                                                                                                                                                                      \\
    \midrule
    \endhead

    % --- PIE INTERMEDIO ---
    \midrule
    \multicolumn{2}{r}{\textit{Continúa en la siguiente página...}}                                                                                                                                               \\
    \endfoot

    % --- PIE FINAL ---
    \bottomrule
    \endlastfoot

    Claridad
                      & Permite verificar si los conceptos, roles (especialmente el Social Guide) y procesos están definidos de forma precisa, con una semántica adecuada que evite ambigüedades en la ejecución. \\

    Completitud
                      & Evalúa si el modelo abarca todas las dimensiones necesarias para la gestión de la deuda social, integrando adecuadamente los eventos adaptados y los artefactos sociales requeridos.      \\

    Idoneidad
                      & Valora la pertinencia de la propuesta frente a las necesidades reales de los equipos de desarrollo y su capacidad para abordar problemas de mejora en entornos de software.               \\

    Facilidad de uso
                      & Determina si los protocolos y herramientas propuestas pueden ser implementados de manera sencilla por el equipo, sin generar una complejidad técnica o administrativa innecesaria.        \\

    Agilidad
                      & Este criterio es el eje central de la validación metodológica y se operacionaliza a través del framework AgilityPRO.                                                                      \\
\end{longtable}
\endgroup

La evaluación de la agilidad en este instrumento no se fundamenta en percepciones subjetivas, sino que tiene como propósito validar la pertinencia y el cumplimiento de los 33 indicadores de agilidad (IA\_01 a IA\_33) definidos en la capa metodológica del framework AgilityPRO \cite{ortega2019agilityevaluation}. Específicamente, las preguntas de esta dimensión buscan que el experto determine si los componentes de SocialScrum ---como el rol del Social Guide, los Eventos Adaptados y los Artefactos Sociales--- operan efectivamente como ``aspectos de agilidad''. Al utilizar el modelo de referencia AgilityRef, esta validación garantiza que SocialScrum no solo gestiona la deuda social, sino que lo hace manteniendo la alineación formal con los valores y principios del Manifiesto Ágil bajo un marco de referencia internacionalmente reconocido.

\subsection{Escala de medición}
Para la evaluación cuantitativa, se utilizó una escala Likert de de 1 a 5 puntos, donde 1 representaba "Totalmente en desacuerdo" y 5 "Totalmente de acuerdo". esto permitió a los expertos expresar su grado de acuerdo o desacuerdo con afirmaciones relacionadas con cada criterio de evaluación. La escala se definió de la siguiente manera en la Tabla~\ref{tab:estructura-instrumento}:

\begingroup
\scriptsize
\renewcommand{\arraystretch}{1.4}
\setlength{\tabcolsep}{3pt}

\begin{longtable}{
    >{\RaggedRight\hspace{0pt}}p{9.2cm}   % Interpretación
    >{\Centering\hspace{0pt}}p{2.2cm}     % Valor
    }

    % --- ENCABEZADO PRIMERA PÁGINA ---
    \caption{Estructura del instrumento de evaluación}
    \label{tab:estructura-instrumento}                                      \\
    \toprule
    \textbf{Interpretación} & \textbf{Valor}                                \\
    \midrule
    \endfirsthead

    % --- ENCABEZADO PÁGINAS SIGUIENTES ---
    \multicolumn{2}{c}{{\bfseries \tablename\ \thetable{} -- Continuación}} \\
    \toprule
    \textbf{Interpretación} & \textbf{Valor}                                \\
    \midrule
    \endhead

    % --- PIE INTERMEDIO ---
    \midrule
    \multicolumn{2}{r}{\textit{Continúa en la siguiente página...}}         \\
    \endfoot

    % --- PIE FINAL ---
    \bottomrule
    \endlastfoot

    Totalmente en desacuerdo
                            & 1                                             \\

    Parcialmente en desacuerdo
                            & 2                                             \\

    Neutralidad
                            & 3                                             \\

    Parcialmente de acuerdo
                            & 4                                             \\

    Totalmente de acuerdo
                            & 5                                             \\
\end{longtable}
\endgroup

\section{Definición del instrumento de validación}
La validación del contenido del marco ágil SocialScrum se organizó en torno a un instrumento diseñado específicamente para este fin, denominado Instrumento de evaluación por juicio de expertos - marco ágil SocialScrum. Este instrumento fue diseñado como un medio para valorar la calidad técnica y la aplicabilidad metodológica del proceso propuesto además de un documento de contextualización entregado a los expertos (ver Anexo~\ref{anexo:contextualizacion}).

\subsection{Estructura de la encuesta}
El instrumento de evaluación del marco ágil SocialScrum se estructuró como un cuestionario en formato digital, compuesto por 20 preguntas cerradas y 3 preguntas abiertas. Las
preguntas cerradas fueron diseñadas para ser valoradas por los expertos en función de cinco criterios de calidad presentados anteriormente, utilizando la escala Likert. Las preguntas abiertas permitieron a los expertos proporcionar comentarios cualitativos y sugerencias de mejora.

\subsection{Preguntas formuladas}
A continuación se presenta el cuestionario dividido en Parte A (preguntas cerradas con escala Likert de 5 puntos, donde 1 = Totalmente en desacuerdo y 5 = Totalmente de acuerdo) y Parte B (preguntas abiertas). Cada pregunta está identificada con un código (P1–P20 para ítems Likert, PA1–PA3 para abiertas):

\begingroup
\scriptsize
\renewcommand{\arraystretch}{1.4}
\setlength{\tabcolsep}{4pt}
\setlength{\LTleft}{0pt}
\setlength{\LTright}{0pt}

\begin{longtable}[H]{
    >{\RaggedRight\hspace{0pt}}p{1.2cm} % Id.
    >{\RaggedRight\hspace{0pt}}p{2.8cm} % Criterio
    >{\RaggedRight\hspace{0pt}}p{\dimexpr\textwidth-1.2cm-2.8cm-6\tabcolsep\relax} % Pregunta
    }

    % --- ENCABEZADO PRIMERA PÁGINA ---
    \caption{Ejemplos de preguntas del instrumento de evaluación}
    \label{tab:ejemplos-preguntas}                                                                                                                                                                                                                                                                                                   \\
    \toprule
    \textbf{Id.} & \textbf{Criterio} & \textbf{Pregunta}                                                                                                                                                                                                                                                                             \\
    \midrule
    \endfirsthead

    % --- ENCABEZADO PÁGINAS SIGUIENTES ---
    \multicolumn{3}{c}{{\bfseries \tablename\ \thetable{} -- Continuación}}                                                                                                                                                                                                                                                          \\
    \toprule
    \textbf{Id.} & \textbf{Criterio} & \textbf{Pregunta}                                                                                                                                                                                                                                                                             \\
    \midrule
    \endhead

    % --- PIE INTERMEDIO ---
    \midrule
    \multicolumn{3}{r}{\textit{Continúa en la siguiente página...}}                                                                                                                                                                                                                                                                  \\
    \endfoot

    % --- PIE FINAL ---
    \bottomrule
    \multicolumn{3}{l}{\footnotesize\textit{Acrónimos utilizados: Identificador (Id.), Pregunta Cerrada (P) y Pregunta Abierta (PA).}}                                                                                                                                                                                               \\
    \endlastfoot

    % =========================
    % PARTE A
    % =========================
    \multicolumn{3}{c}{\textbf{Parte A – Preguntas Cerradas}}                                                                                                                                                                                                                                                                        \\* \nopagebreak
    \midrule

    P1           & Claridad          & ¿Considera que los conceptos del modelo SocialScrum están claramente definidos?                                                                                                                                                                                                               \\
    P2           & Claridad          & ¿El rol del Social Guide está claramente diferenciado de otros roles (por ejemplo, con el Scrum Master, Product Owner o el departamento de RRHH) evitando conflictos de autoridad?                                                                                                            \\
    P3           & Claridad          & ¿El modelo establece mecanismos claros de gobernanza para resolver desacuerdos entre el Social Guide y otros roles (por ejemplo: protocolo de escalamiento PO-Social Guide, mediación del Scrum Master o arbitraje del CTO), evitando impactos negativos en la toma de decisiones del equipo? \\
    \midrule
    P4           & Completitud       & ¿Cree que el modelo SocialScrum cubre todos los aspectos necesarios para la gestión ágil de proyectos?                                                                                                                                                                                        \\
    P5           & Completitud       & ¿Los artefactos sociales introducidos (por ejemplo, el Social Product Backlog, las historias de usuario sociales y el registro de community smells) están bien definidos y facilitan el seguimiento de la salud del equipo?                                                                   \\
    P6           & Completitud       & ¿El modelo incluye mecanismos de auditoría y transparencia que permiten al equipo verificar la actuación del Social Guide, previniendo posibles abusos de autoridad o desviaciones éticas?                                                                                                    \\
    \midrule
    P7           & Idoneidad         & ¿Los valores fundamentales de Scrum (coraje, foco, compromiso, apertura y respeto) se reflejan adecuadamente en las prácticas de SocialScrum para manejar la deuda social del equipo?                                                                                                         \\
    P8           & Idoneidad         & ¿La métrica Health Score proporciona una evaluación objetiva y útil de la salud social del equipo, sirviendo de base para tomar acciones de mejora cuando sea necesario?                                                                                                                      \\
    P9           & Idoneidad         & ¿SocialScrum garantiza la confidencialidad de la información personal compartida por el equipo, de modo que ningún dato sensible se comunica a terceros (por ejemplo, gerencia o RRHH) sin consentimiento expreso?                                                                            \\
    \midrule
    P10          & Facilidad de Uso  & ¿El perfil y las competencias del Social Guide están claramente definidos y son viables, complementando los roles existentes en el equipo sin solapamiento de responsabilidades?                                                                                                              \\
    P11          & Facilidad de Uso  & ¿La implementación de SocialScrum en equipos ágiles reales es factible y no requiere cambios organizacionales o culturales impracticables?                                                                                                                                                    \\
    P12          & Facilidad de Uso  & ¿La sobrecarga operacional de SocialScrum es sostenible para equipos ágiles a mediano/largo plazo?                                                                                                                                                                                            \\
    \midrule
    P13          & Agilidad          & ¿La adaptación de eventos Scrum en SocialScrum (como el Social Daily Standup o las retrospectivas sociales) se integra de forma fluida y aporta un valor operativo que justifica su implementación?                                                                                           \\
    P14          & Agilidad          & ¿El sistema de estimación Social Story Points (SSP) es compatible con las prácticas ágiles estándar y resulta comprensible para todos los miembros del equipo?                                                                                                                                \\
    P15          & Agilidad          & ¿El mecanismo de priorización socio-técnica de SocialScrum (integrando intervenciones sociales junto con tareas técnicas) es práctico y logra equilibrar la mejora del equipo con la entrega de valor al producto?                                                                            \\

    % =========================
    % PARTE B
    % =========================
    \midrule
    \multicolumn{3}{c}{\textbf{Parte B – Preguntas Abiertas}}                                                                                                                                                                                                                                                                        \\* \nopagebreak
    \midrule

    PA1          & ---               & ¿Qué elementos del modelo SocialScrum considera usted que faltan o están poco desarrollados?                                                                                                                                                                                                  \\
    PA2          & ---               & ¿Qué sugerencias de mejora específicas propondría para el modelo SocialScrum? Justifique su respuesta con fundamentos teóricos o prácticos.                                                                                                                                                   \\
\end{longtable}
\endgroup


\section{Perfil y selección de expertos participantes}
\subsection{Perfilamiento de los participantes}
Para la validación del modelo SocialScrum, se seleccionaron expertos con experiencia en metodologías ágiles, desarrollo de software y gestión de proyectos. Los criterios de selección incluyeron:
\begin{itemize}
    \item Experiencia académica o profesional en el área de desarrollo de software.
    \item Experiencia mínima de 2 años en la aplicación de metodologías ágiles.
    \item Participación en proyectos ágiles exitosos.
    \item Conocimiento en herramientas de colaboración y gestión de proyectos ágiles.
    \item Preferiblemente, experiencia en equipos distribuidos o remotos.
\end{itemize}

\subsection{Selección de participantes}
Se contactaron a 8 expertos potenciales a través de redes profesionales y académicas. De estos, los 8 aceptaron participar en la validación del modelo SocialScrum. Los participantes provinieron de diversas industrias, incluyendo desarrollo de software, consultoría ágil y educación en ingeniería de software. A continuación, se presenta un resumen del perfil de los expertos participantes en la validación del modelo SocialScrum (ver Tabla~\ref{tab:perfil-expertos}).

\begingroup
\scriptsize
\renewcommand{\arraystretch}{1.4}
\setlength{\tabcolsep}{3pt}

% Columna de texto alineada a la izquierda con salto de línea
\newcolumntype{Y}{>{\raggedright\arraybackslash}p}
% Columna centrada SOLO para datos cortos
\newcolumntype{C}{>{\centering\arraybackslash}p}

\begin{longtable}{
        C{1.2cm}   % Id.
        Y{3.4cm}   % Último título
        C{2.4cm}   % Años de experiencia
        Y{4.4cm}   % Cargo actual
        Y{3.6cm}   % Metodologías
    }

    % --- ENCABEZADO PRIMERA PÁGINA ---
    \caption{Perfil de los expertos participantes en la validación del modelo SocialScrum}
    \label{tab:perfil-expertos}                                                                                                                                    \\
    \toprule
    \textbf{Id.} & \textbf{Último título obtenido} & \textbf{Años de experiencia} & \textbf{Cargo actual}                 & \textbf{Metodologías ágiles conocidas} \\
    \midrule
    \endfirsthead

    % --- ENCABEZADO PÁGINAS SIGUIENTES ---
    \multicolumn{5}{c}{{\bfseries \tablename\ \thetable{} -- Continuación}}                                                                                        \\
    \toprule
    \textbf{Id.} & \textbf{Último título obtenido} & \textbf{Años de experiencia} & \textbf{Cargo actual}                 & \textbf{Metodologías ágiles conocidas} \\
    \midrule
    \endhead

    % --- PIE INTERMEDIO ---
    \midrule
    \multicolumn{5}{r}{\textit{Continúa en la siguiente página...}}                                                                                                \\
    \endfoot

    % --- PIE FINAL ---
    \bottomrule
    \multicolumn{5}{l}{\footnotesize \textit{Acrónimos utilizados: Identificador (Id.) y Experto (E).}}                                                            \\
    \endlastfoot

    E1           & Ingeniero de sistemas           & 4 años                       & Desarrollador Senior de software      & Scrum, Kanban, XP                      \\
    E2           & Ingeniero de sistemas           & 4 años                       & Desarrollador Semi-Senior de software & Scrum, Lean, Kanban                    \\
    E3           & Ingeniero de sistemas           & 9 años                       & Desarrollador Senior de software      & Scrum, Kanban                          \\
    E4           & Ingeniero de sistemas           & 5 años                       & Desarrollador Senior de software      & Scrum, XP, Lean                        \\
    E5           & Ingeniero de sistemas           & 4 años                       & Desarrollador Semi-Senior de software & Scrum, Kanban                          \\
    E6           & Ingeniero en Telemática         & 5 años                       & Ingeniero DevOps Senior               & Scrum, Kanban, SAFe                    \\
    E7           & Ingeniero de sistemas           & 3 años                       & Desarrollador Semi-Senior de software & Scrum, Lean                            \\
    E8           & Ingeniero de sistemas           & 5 años                       & Desarrollador Senior de software      & Scrum, Kanban, XP, Lean                \\
\end{longtable}
\endgroup

\section{Protocolo de validación mediante juicio de expertos}
Para la evaluación del modelo SocialScrum, se diseñó un protocolo de validación basado en la técnica de juicio de expertos. Para la evaluación y validación de este modelo se incluyó un documento entregado a los participantes con la contextualización del tema a evaluar (ver Anexo~\ref{anexo:contextualizacion}). Esta técnica permite obtener una valoración cualitativa y cuantitativa del modelo a partir de la experiencia y conocimientos de profesionales en metodologías ágiles y desarrollo de software. El protocolo se estructuró en varias etapas, que se describen a continuación.

\subsection{Objetivo de la validación}
El objetivo principal de esta validación es evaluar la adecuación del modelo SocialScrum en términos de su alineación con los principios ágiles, su aplicabilidad en contextos reales de desarrollo de software y su capacidad para fomentar la colaboración y eficiencia en equipos distribuidos.

\section{Resultados de la validación y análisis cuantitativo}
En esta sección se presentan los resultados obtenidos a partir de la validación del modelo SocialScrum. El propósito central de este análisis es determinar en qué medida el proceso cumple con los criterios de calidad previamente definidos —claridad, completitud, idoneidad, facilidad de uso y agilidad—, así como identificar percepciones cualitativas
que permitan reconocer sus fortalezas y oportunidades de mejora

\subsection{Análisis cuantitativo de las respuestas}
Se recopilaron las respuestas de los 8 expertos participantes, quienes evaluaron cada una de las 20 preguntas del instrumento de validación utilizando la escala Likert propuesta. A continuación, se presentan los resultados cuantitativos obtenidos para cada criterio de evaluación.

\begin{table}[!htbp]
    \caption{Distribución de respuestas de los expertos por pregunta en la escala de Likert.}
    \label{tab:dist-resp-likert}
    \centering
    \small
    \begin{tabularx}{\textwidth}{|c|Y|c|c|c|c|c|}
        \hline
        \multirow{2}{*}{\textbf{Id.}}                           &
        \multirow{2}{*}{\textbf{Preguntas con escala discreta}} &
        \multicolumn{5}{c|}{\textbf{Número de respuestas por nivel de la escala}}                                                                                                                                                                                               \\
        \cline{3-7}
                                                                &                                                                                                           & \textbf{Escala 5} & \textbf{Escala 4} & \textbf{Escala 3} & \textbf{Escala 2} & \textbf{Escala 1} \\
        \hline
        P1.                                                     & ¿El objetivo general del marco ágil SocialScrum está claramente formulado?                                & 15                & 4                 & 1                 & 0                 & 0                 \\
        \hline
        P2.                                                     & ¿Las actividades descritas son entendibles por un lector con conocimientos básicos en mejora de procesos? & 14                & 5                 & 1                 & 0                 & 0                 \\
        \hline
        P3.                                                     & ¿La descripción de roles, entradas, salidas y productos es clara y accesible?                             & 13                & 6                 & 1                 & 0                 & 0                 \\
        \hline
        P4.                                                     & ¿El lenguaje utilizado evita ambigüedades o tecnicismos innecesarios?                                     & 12                & 7                 & 1                 & 0                 & 0                 \\
        \hline
    \end{tabularx}
    \smallskip
    \raggedright
    \textit{Acrónimos utilizados: Identificador (Id.).}
\end{table}

A continuación, se presentan los resultados organizados por cada uno de los criterios de evaluación considerados en el instrumento: (i) claridad, (ii) completitud, (iii) idoneidad,
(iv) facilidad de uso y (v) coherencia. En las subsecciones siguientes se muestran las tablas y figuras que resumen el comportamiento de las respuestas para cada criterio, lo cual permite disponer de una visión diferenciada del nivel de conformidad alcanzado en cada dimensión del modelo SocialScrum.

\subsubsection{Claridad}
\begin{table}[!htbp]
    \caption{Distribución de respuestas de los expertos para el criterio de claridad.}
    \label{tab:dist-resp-claridad}
    \centering
    \begin{tabularx}{\textwidth}{|c|Y|c|c|c|c|c|}
        \hline
        \textbf{Id.} & \textbf{Pregunta}                                                               & \textbf{Escala 5}                                             & \textbf{Escala 4} & \textbf{Escala 3} & \textbf{Escala 2} & \textbf{Escala 1} \\
        \hline
        P1           & ¿Considera que los conceptos del modelo SocialScrum están claramente definidos? & [Pendiente de consolidación tras la recepción de formularios] & 0                 & 0                 & 0                 & 0                 \\
        \hline
    \end{tabularx}
    \smallskip
    \raggedright
    \textit{Acrónimos utilizados: Identificador (Id.).}
\end{table}

\subsubsection{Completitud}
\begin{table}[!htbp]
    \caption{Distribución de respuestas de los expertos para el criterio de completitud.}
    \label{tab:dist-resp-completitud}
    \centering
    \begin{tabularx}{\textwidth}{|c|Y|c|c|c|c|c|}
        \hline
        \textbf{Id.} & \textbf{Pregunta}                                                               & \textbf{Escala 5}                                             & \textbf{Escala 4} & \textbf{Escala 3} & \textbf{Escala 2} & \textbf{Escala 1} \\
        \hline
        P1           & ¿Considera que los conceptos del modelo SocialScrum están claramente definidos? & [Pendiente de consolidación tras la recepción de formularios] & 0                 & 0                 & 0                 & 0                 \\
        \hline
    \end{tabularx}
    \smallskip
    \raggedright
    \textit{Acrónimos utilizados: Identificador (Id.).}
\end{table}

\subsubsection{Idoneidad}
\begin{table}[!htbp]
    \caption{Distribución de respuestas de los expertos para el criterio de idoneidad.}
    \label{tab:dist-resp-idoneidad}
    \centering
    \begin{tabularx}{\textwidth}{|c|Y|c|c|c|c|c|}
        \hline
        \textbf{Id.} & \textbf{Pregunta}                                                               & \textbf{Escala 5}                                             & \textbf{Escala 4} & \textbf{Escala 3} & \textbf{Escala 2} & \textbf{Escala 1} \\
        \hline
        P1           & ¿Considera que los conceptos del modelo SocialScrum están claramente definidos? & [Pendiente de consolidación tras la recepción de formularios] & 0                 & 0                 & 0                 & 0                 \\
        \hline
    \end{tabularx}
    \smallskip
    \raggedright
    \textit{Acrónimos utilizados: Identificador (Id.).}
\end{table}

\subsubsection{Facilidad de uso}
\begin{table}[!htbp]
    \caption{Distribución de respuestas de los expertos para el criterio de facilidad de uso.}
    \label{tab:dist-resp-facilidad-uso}
    \centering
    \begin{tabularx}{\textwidth}{|c|Y|c|c|c|c|c|}
        \hline
        \textbf{Id.} & \textbf{Pregunta}                                                               & \textbf{Escala 5}                                             & \textbf{Escala 4} & \textbf{Escala 3} & \textbf{Escala 2} & \textbf{Escala 1} \\
        \hline
        P1           & ¿Considera que los conceptos del modelo SocialScrum están claramente definidos? & [Pendiente de consolidación tras la recepción de formularios] & 0                 & 0                 & 0                 & 0                 \\
        \hline
    \end{tabularx}
    \smallskip
    \raggedright
    \textit{Acrónimos utilizados: Identificador (Id.).}
\end{table}

\subsubsection{Agilidad}
\begin{table}[!htbp]
    \caption{Distribución de respuestas de los expertos para el criterio de agilidad.}
    \label{tab:dist-resp-agilidad}
    \centering
    \begin{tabularx}{\textwidth}{|c|Y|c|c|c|c|c|}
        \hline
        \textbf{Id.} & \textbf{Pregunta}                                                               & \textbf{Escala 5}                                             & \textbf{Escala 4} & \textbf{Escala 3} & \textbf{Escala 2} & \textbf{Escala 1} \\
        \hline
        P1           & ¿Considera que los conceptos del modelo SocialScrum están claramente definidos? & [Pendiente de consolidación tras la recepción de formularios] & 0                 & 0                 & 0                 & 0                 \\
        \hline
    \end{tabularx}
    \smallskip
    \raggedright
    \textit{Acrónimos utilizados: Identificador (Id.).}
\end{table}

\subsection{Resultados preguntas abiertas}
El presente análisis se enfoca en las respuestas abiertas proporcionadas por los expertos durante la validación del modelo SocialScrum. Estas respuestas cualitativas ofrecen una perspectiva enriquecedora sobre las fortalezas y áreas de mejora del marco ágil propuesto. A continuación, se presentan las observaciones y recomendaciones más relevantes extraídas de las respuestas de los expertos

\begin{table}
    \centering
    \begin{tabularx}{\textwidth}{|c|Y|c|c|}
        \hline
        \textbf{Id.} & \textbf{Observación / Recomendación}                                                                                                      & \textbf{Calificación} & \textbf{Justificación}                                        \\
        \hline
        PA1          & El modelo SocialScrum presenta una estructura clara y bien definida, lo que facilita su comprensión y aplicación en equipos distribuidos. & Si                    & [Pendiente de consolidación tras la recepción de formularios] \\
        \hline
        PA2          & Se recomienda incluir ejemplos prácticos o casos de estudio que ilustren la implementación del modelo en diferentes contextos.            & Si                    & [Pendiente de consolidación tras la recepción de formularios] \\
        \hline
        PA3          & Sería beneficioso proporcionar guías adicionales sobre cómo adaptar el modelo a las particularidades de cada equipo o proyecto.           & No                    & [Pendiente de consolidación tras la recepción de formularios] \\
        \hline
    \end{tabularx}
    \smallskip
    \raggedright
    \textit{Acrónimos utilizados: Identificador (Id.).}
    \caption{Observaciones y recomendaciones de los expertos}
    \label{tab:observaciones-recomendaciones}
\end{table}

\section{Desempeño del modelo y recomendaciones de mejora}

\subsubsection{Fortalezas del modelo SocialScrum}
\begin{itemize}
    \item \textbf{Claridad y estructura}: Los expertos destacaron que el modelo SocialScrum presenta una estructura clara y bien definida, lo que facilita su comprensión y aplicación en equipos distribuidos.
    \item \textbf{Alineación con principios ágiles}: Se reconoció que el modelo está alineado con los valores y principios del Manifiesto Ágil, promoviendo la colaboración y la adaptabilidad.
    \item \textbf{Enfoque en la colaboración}: Los expertos valoraron positivamente el énfasis del modelo en fomentar la colaboración entre los miembros del equipo, especialmente en contextos distribuidos.
\end{itemize}

\subsubsection{Debilidades del modelo SocialScrum}
\begin{itemize}
    \item \textbf{Falta de ejemplos prácticos}: Algunos expertos señalaron la ausencia de ejemplos prácticos o casos de estudio que ilustren la implementación del modelo en diferentes contextos.
    \item \textbf{Adaptabilidad limitada}: Se mencionó que el modelo podría beneficiarse de guías adicionales sobre cómo adaptarlo a las particularidades de cada equipo o proyecto.
    \item \textbf{Clarificación de roles}: Algunos expertos indicaron que ciertos roles y responsabilidades dentro del modelo podrían no estar suficientemente claros, lo que podría generar confusiones durante su implementación.
    \item \textbf{Falta de herramientas de soporte}: Se observó que el modelo no proporciona recomendaciones sobre herramientas tecnológicas que puedan facilitar su aplicación en equipos distribuidos.
    \item \textbf{Complejidad en la implementación}: Algunos expertos consideraron que la implementación del modelo podría resultar compleja para equipos con poca experiencia en metodologías ágiles.
\end{itemize}

\subsubsection{Oportunidades de mejora}
\begin{itemize}
    \item \textbf{Incorporación de ejemplos prácticos}: Se recomendó incluir ejemplos prácticos o casos de estudio que ilustren la implementación del modelo en diferentes contextos.
    \item \textbf{Guías de adaptación}: Los expertos sugirieron proporcionar guías adicionales sobre cómo adaptar el modelo a las particularidades de cada equipo o proyecto.
    \item \textbf{Clarificación de roles}: Algunos expertos señalaron la necesidad de clarificar los roles y responsabilidades dentro del modelo para evitar confusiones durante su implementación.
    \item \textbf{Herramientas de soporte}: Se propuso la inclusión de recomendaciones sobre herramientas tecnológicas que puedan facilitar la aplicación del modelo en equipos distribuidos.
\end{itemize}


